\documentclass[11pt,a4paper]{article}
\usepackage[margin=23mm]{geometry}
\usepackage{amsmath,amssymb,amsthm}
\setlength{\emergencystretch}{2em}
\newtheorem{theorem}{Theorem}
\newcommand{\gen}[1]{\langle#1\rangle}
\newcommand{\Alt}{\operatorname{Alt}}
\title{The last fixed-point budgets at odd degree}
\author{Human Compression project: L044 continuation}
\date{}
\begin{document}\maketitle
Let $G=\gen{\rho_0,\ldots,\rho_{r-1}}\le\Alt(\Omega)$
be transitive with a connected string C-marking and all
maximal omissions intransitive. Assume $r\ge6$, $n=2r+1$
and the actual transitive six-point head on $D=\{1,\ldots,6\}$,
with the literal interface
\[
 \tau=\rho_3=(6\ b)\delta,\quad
 R=G_{\ge3}=\Alt(\Omega\setminus\{1,\ldots,5\}),\quad
 \sigma=\rho_4,\quad \sigma|_D=1.
\]
The odd involution $\delta$ fixes $D,b$. Put $H=G_{\ge5}$
and $K=G_{>5}$, with $H$ fixing $D,b$. Suppose each tail
colour has two full ambient-eligible edges in one actual
$H$-orbit, and normalize these assignments by the acquired
whole-orbit transport.

\begin{theorem}
Under these hypotheses $H$ cannot have nine or more fixed points.
\end{theorem}
\begin{proof}
Write $f=|\operatorname{Fix}(H)|$, and let $c\ge1$ count
its nontrivial orbits. The complete ambient interval theorem
applies to all $r-5$ tail colours and gives
\[
                  2(r-5)\le n-f-c,\qquad f+c\le11.       \tag{1}
\]
Thus $f\ge9$ leaves only $f=9,c\le2$ or $f=10,c=1$.
The acquired actual head argument gives
\[
             x=b^\sigma\ne b,\qquad \delta x=x,\qquad Kx=x.
                                                               \tag{2}
\]
Moreover $x$ is not fixed by $H$. If it were, $H$ would
fix $D\cup\{b,x\}$ pointwise. The prefix $G_{\le4}$
preserves that set: the head preserves $D$, $\tau$ fixes
$x$ and exchanges $6,b$, and $\sigma$ fixes $D$ and
exchanges $b,x$. The full group would preserve its
eight points, contrary to transitivity and $n\ge13$.

Suppose first that $c=1$. The involution $\delta$ preserves
the unique nontrivial $H$-orbit and fixes its point $x$,
so centralization makes it trivial on that orbit.
Outside $D,b$ there are at most three $H$-fixed points.
As $\delta$ is odd, it is literally one transposition
on those points. The complete acquired exterior-pair
obstruction applies with every required ambient eligible
pair, and excludes this case.

We may therefore assume $f=9,c=2$. For the two orbits put
$d_i=|F_i|-|O_i|/2$. The normalized bounds give
\[
       d_1+d_2=-1,\qquad d_p\le0,
       \qquad d_i<0\quad(i\ne p),                      \tag{3}
\]
where $5\in F_p$. The involution $\delta$ preserves
both orbits. Otherwise they have equal size $m$;
transporting their assigned pairs by $\tau$ and using
the selected-colour bound gives combined deficit at
most $-\lceil m/2\rceil\le-2$ for $m\ge3$.
For $m=2$ both assignments are empty and the deficit
is exactly $-2$. Either conclusion contradicts (3).

If $\delta$ fixes both orbits pointwise, its odd action
is again the single transposition of the two additional
fixed points, already excluded. Otherwise let $W$ be
an orbit with nontrivial internal action. It is free
there and has even size $2q$. Its deficit is at most
$-1$: use the complete pair-forest bound when its
assignment is nonempty, and its exact deficit otherwise.
It cannot be the first assigned orbit, since the other
deficit is negative. Equation (3) therefore forces
\[
        d_W=-1,\qquad d_V=0,\qquad V=O_p.
\]
The involution $\delta$ fixes $V$ pointwise; otherwise
the same bound would make $d_V\le-1$. The point $x$
lies in $V$, by (2) and the preceding nonfixation
argument. Every possible further action of $\delta$
is on the two additional $H$-fixed points.

If $q=1$, its assignment is empty, and oddness forces
$\delta$ to fix those two extra points. It is the
literal transposition on $W$, again impossible.
If $q\ge3$, then $F_W$ has $q-1$ colours, all at
least six. Its normalized $K$-action is intransitive.
The complete odd-degree row theorem gives a last
assigned label $t<r-1$ with
$|\operatorname{Fix}(G_{>t})|\ge2t+3$.
The remaining nonempty paired interval has support
at least $2(r-1-t)+1$, forcing $n\ge2r+2$.
Thus only $q=2$ remains.

In this case $|W|=4$, $F_W=\{j\}$ with $j\ge6$,
and $F_V$ is the complementary set of tail colours.
As $|V|=2|F_V|$ and $K$ fixes $x\in V$, the
complete fixed-anchor lemma makes $F_V$ an initial
interval. Its unique missing label must therefore
be $j=r-1$. Put $Q=\gen{\rho_j}$.

We give the exact two-row argument with the present
fixed points. The selected eligible $j$-edge joins
distinct $\delta$-pairs in $W$. No other tail mark
can exchange those pair vertices: it and $\tau$
would then connect all of $W$ in $G_j$, contrary
to eligibility. Thus $K$ is transitive on the two
pair vertices. Normalization makes it intransitive
on $W$, so it has two orbits of size two, each
meeting each pair once. Its pair-trivial elements
fix both rows pointwise. Consequently $\rho_j|_W$
is a double transposition and every other mark
of $K$ fixes $W$ pointwise.

For $4\le l\le j-2$ put $Z_l=W^{G_{\le l}}$.
Each is a union of nontrivial two-point $Q$-orbits,
since the prefix centralizes $Q$. We have $|Z_4|\ge6$:
otherwise $G_{\le4}$ and $H$ would both preserve $W$,
contradicting transitivity. Every $Z_l$ is disjoint
from the ten $Q$-fixed points
\[
                         \operatorname{Fix}(H)\cup\{x\}.
\]
For $4\le l\le j-3$, every mark of index at least
$l+2$ except $\rho_j$ fixes $W$ and commutes with
the prefix, so fixes $Z_l$. The group $Q$ and the
prefix preserve $Z_l$. Transitivity forces
$\rho_{l+1}$ to enlarge this proper set, by at
least one whole two-point $Q$-orbit. Therefore
\[
        |Z_{j-2}|\ge6+2(j-6)=2j-6,
        \qquad n\ge(2j-6)+10=2j+4.
\]
But $j=r-1$ gives $n=2j+3$, the final contradiction.
\end{proof}

\paragraph{Attribution and boundary.}
The complete interval bound, normalized deficit theorem,
fixed-anchor theorem, odd-degree row exit and literal
exterior-pair obstruction are acquired. The last paragraph
adapts the acquired four-point-row proof with its ten
actual fixed points stated explicitly. The new step is
the exhaustive remaining fixed-point budget. Combined
with an independently reviewed lower bound $f\ge9$,
this would exclude the whole single-orbit assignment
configuration; that combination is not a premise here.
Split assignments and their source routing remain.
No new induction, computation, formal, original or
question credit is claimed.
\end{document}
